\chapter{Storage, Indexing, and Search}
This chapter describes the persistence layer and retrieval mechanisms that power the internship search platform. We detail the dual-storage architecture designed to support both structured metadata queries and high-dimensional semantic search, followed by the hybrid retrieval strategy that combines lexical and dense vector scores.

\section*{Chapter Outline}
This chapter is organized as follows:
\begin{itemize}
    \item \textbf{Section 9.1: Data Storage Model} describes the dual-storage architecture using PostgreSQL and pgvector.
    \item \textbf{Section 9.2: Search Index Design} details the lexical (BM25) and semantic (HNSW) indexing strategies.
    \item \textbf{Section 9.3: Hybrid Retrieval Strategy} explains the score fusion using Reciprocal Rank Fusion (RRF).
    \item \textbf{Section 9.4: API Endpoints} outlines the RESTful interfaces for search and data ingestion.
\end{itemize}

\section{Data Storage Model}
The storage layer is designed to balance the need for relational integrity with the performance requirements of vector-based search. We adopt a \textbf{dual-storage architecture}.

\subsection{Relational Storage (PostgreSQL)}
PostgreSQL serves as the system of record for all persistent data. It handles:
\begin{itemize}
    \item \textbf{Core Metadata:} Storing the `StructuredDocument` fields defined in Chapter 7, including job titles, company names, and canonical IDs for occupations and skills.
    \item \textbf{System State:} Managing ingestion logs, source tracking (URLs, file hashes), and user-generated data (e.g., saved searches).
    \item \textbf{JSONB Integration:} We utilize the `JSONB` data type for semi-structured fields that may vary between sources, such as specific "benefits" or "nice-to-have" requirements that are not yet part of the canonical schema.
\end{itemize}

\subsection{Vector Storage (pgvector)}
For semantic retrieval, we utilize the `pgvector` extension for PostgreSQL. This allows us to store the 384-dimensional SBERT embeddings directly alongside the relational data. By keeping embeddings in the same database, we simplify the architecture and avoid the synchronization challenges associated with managing a separate vector database (like Milvus or Pinecone).

\section{Search Index Design}
To provide low-latency search results, we maintain two distinct types of indices: lexical and semantic.

\subsection{Lexical Index (GIN/BM25)}
A Generalized Inverted Index (GIN) is built on the cleaned description and title fields.
\begin{itemize}
    \item \textbf{Tokenization:} Text is tokenized using a domain-specific analyzer that preserves technical abbreviations (e.g., "C\#", ".NET").
    \item \textbf{Ranking Model:} We configure the text search to use a BM25-based scoring algorithm, which provides superior relevance for exact keyword matches compared to standard TF-IDF.
\end{itemize}

\subsection{Semantic Index (HNSW)}
The dense vectors are indexed using the \textbf{Hierarchical Navigable Small World (HNSW)} algorithm.
\begin{itemize}
    \item \textbf{Parameters:} We set the construction parameters to $M=16$ and $ef\_construction=64$, providing a balance between indexing speed and retrieval recall.
    \item \textbf{Approximate Search:} HNSW allows the system to find the "Top-K" most semantically similar documents in sub-linear time, which is essential for real-time user interaction.
\end{itemize}

\section{Hybrid Retrieval Strategy}
The core search functionality employs a hybrid strategy to leverage the precision of keyword matching and the recall of semantic search.

\subsection{Reciprocal Rank Fusion (RRF)}
To combine scores from the BM25 index and the HNSW vector index, we implement \textbf{Reciprocal Rank Fusion (RRF)}. As defined in Equation \ref{eq:rrf}, RRF allows us to merge two ranked lists without needing to normalize the raw scores from disparate systems. This ensures that a document that ranks high in \textit{either} lexical or semantic search remains highly visible to the user.

\subsection{The Retrieval Pipeline}
The hybrid retrieval follows a two-stage process:
\begin{enumerate}
    \item \textbf{Candidate Retrieval:} The system fetches the top 100 results from the lexical index and the top 100 results from the vector index.
    \item \textbf{Re-ranking and Fusion:} The RRF algorithm merges these lists. Any hard filters (e.g., "Location = New York" or "Skill = Python") are applied as SQL constraints during the initial retrieval to ensure the search remains efficient.
\end{enumerate}

\section{API Endpoints}
The system's functionality is exposed via a RESTful API, facilitating integration with the frontend application and external tools.

\subsection{Core Search API}
The primary endpoint is `GET /api/v1/search`. It accepts the following parameters:
\begin{itemize}
    \item `q`: The natural language query string.
    \item `filters`: A JSON object for hard constraints (skills, level, location).
    \item `alpha`: A weighting parameter (0.0 to 1.0) to manually adjust the balance between lexical and semantic search (default is 0.5).
    \item `limit/offset`: Standard pagination parameters.
\end{itemize}

\subsection{Ingestion and Analytics API}
\begin{itemize}
    \item `POST /api/v1/upload/json`: Handles manual file uploads, triggering the mapping heuristics defined in Section 6.3.
    \item `GET /api/v1/stats`: Returns aggregate data for the dashboard, such as "Top 10 most demanded skills this week" or "Distribution of roles by industry."
\end{itemize}
